
\newgeometry{top=12pt, right=1.75in}
\begin{abstract}


Traffic congestion significantly impacts growing urban centers like Iloilo City. Expanding road capacity is an intuitive solution, but adding links can counterintuitively increase overall travel costs, a phenomenon known as the Braess Paradox. This study develops a computational network model to identify the presence of the Braess Paradox within the primary arterial road networks of Iloilo City, using data from the DPWH Atlas 2024. By employing macroscopic static traffic assignment and maximizing the Price of Anarchy (PoA), various network topologies and road closures were simulated. Results confirm the paradox's existence, showing that selfish routing behavior can increase Total System Travel Time by 58.15\%. The model identified 4,627 unidirectional and 3,662 bidirectional link closures that mathematically improved travel time, indicating that congestion is heavily driven by directional imbalances rather than insufficient road space.

\begin{flushleft}
\begin{tabular}{lp{4.3in}}

\hspace{-0.5em}\textbf{Keywords:}\hspace{0.25em} & Braess Paradox, Traffic Congestion, Computational Modeling, Price of Anarchy (PoA), User Equilibrium, Iloilo City, Macroscopic Static Traffic Assignment

\\

\end{tabular}
\end{flushleft}
\end{abstract}
\restoregeometry